\documentclass[a4paper,12pt]{article}

\usepackage[T2A]{fontenc}
\usepackage[utf8]{inputenc}
\usepackage[russian]{babel}
\usepackage{amsmath,amssymb,amsthm,mathtools}
\usepackage{helvet}
\renewcommand{\familydefault}{\sfdefault}
\usepackage{icomma}
\usepackage[a4paper,margin=2.15cm]{geometry}
\usepackage{microtype}
\usepackage{tikz}
\usetikzlibrary{calc}
\usetikzlibrary{arrows.meta,positioning,shapes.geometric}
\usepackage{tkz-euclide}
\usepackage{pgfplots}
\pgfplotsset{compat=1.18}
\usepackage{tabularx,booktabs,longtable}
\usepackage{enumitem}
\usepackage{hyperref}
\usepackage{parskip}
\usepackage{fancyhdr}
\usepackage{setspace}
\usepackage{xcolor}

\definecolor{accent}{HTML}{008080}
\definecolor{darkaccent}{HTML}{004D4D}
\definecolor{textgray}{HTML}{333333}
\definecolor{softgray}{HTML}{F3F5F5}

\hypersetup{
    colorlinks=true,
    linkcolor=darkaccent,
    urlcolor=darkaccent
}

\onehalfspacing
\setlength{\parindent}{0pt}
\setlist[itemize]{
    leftmargin=1.5em,
    itemsep=0.35em,
    topsep=0.35em
}
\setlist[enumerate]{
    leftmargin=1.8em,
    itemsep=0.45em,
    topsep=0.4em
}
\renewcommand{\arraystretch}{1.35}

\pagestyle{fancy}
\fancyhf{}
\fancyhead[L]{\small\textcolor{textgray}{Анна Трапезникова \textbullet{} 10 класс}}
\fancyhead[R]{\small\textcolor{textgray}{ОДЗ и метод интервалов}}
\fancyfoot[C]{\small\thepage}
\renewcommand{\headrulewidth}{0.3pt}
\setlength{\headheight}{14pt}

\newcommand{\topicsection}[1]{%
    \section*{\textcolor{darkaccent}{#1}}%
    \vspace{-0.4em}%
    \textcolor{accent}{\rule{\linewidth}{0.7pt}}%
    \vspace{0.4em}%
}

\newcommand{\key}[1]{%
    \textcolor{darkaccent}{\textbf{#1}}%
}

\tikzset{
    flowstart/.style={
        ellipse,
        draw=darkaccent,
        very thick,
        minimum width=4.2cm,
        minimum height=1.1cm,
        align=center,
        inner sep=5pt
    },
    flowaction/.style={
        rounded corners=4pt,
        draw=accent,
        thick,
        minimum width=10.6cm,
        minimum height=1.15cm,
        text width=9.8cm,
        align=center,
        inner sep=6pt
    },
    flowbranch/.style={
        rounded corners=4pt,
        draw=accent,
        thick,
        minimum width=5.15cm,
        minimum height=1.2cm,
        text width=4.55cm,
        align=center,
        inner sep=5pt
    },
    flowdecision/.style={
        diamond,
        aspect=2.2,
        draw=darkaccent,
        thick,
        text width=5.7cm,
        align=center,
        inner sep=2pt
    },
    flowarrow/.style={
        -{Latex[length=3mm]},
        thick,
        draw=textgray
    },
    flowlabel/.style={
        font=\small,
        fill=white,
        inner sep=1.5pt
    }
}

\begin{document}

% =========================================================
% Титульный лист
% =========================================================

\begin{titlepage}
\thispagestyle{empty}
\centering

\vspace*{2.4cm}

{\Large\textcolor{darkaccent}{\textbf{Чек-лист}}\par}

\vspace{0.7cm}

{\huge\textbf{Область определения функции и метод интервалов}\par}

\vspace{0.55cm}

{\Large Квадратные и дробно-рациональные неравенства\par}

\vspace{2.2cm}

{\large Анна Трапезникова, 10 класс\par}

\vspace{0.8cm}

{\large \today\par}

\vfill

{\large
\textcolor{textgray}{
Лёвин Артём Александрович эксклюзивно для Анны Трапезниковой
}\par}

\vspace{1.2cm}

\begin{tikzpicture}[remember picture,overlay]
\draw[
    accent,
    line width=1.2pt
]
    ([xshift=1.7cm,yshift=1.55cm]current page.south west)
    .. controls
    ([xshift=6.3cm,yshift=2.25cm]current page.south west)
    and
    ([xshift=-6.1cm,yshift=0.9cm]current page.south east)
    ..
    ([xshift=-1.7cm,yshift=1.7cm]current page.south east);
\end{tikzpicture}

\end{titlepage}

% =========================================================
% Основной чек-лист
% =========================================================

\topicsection{1. Что нужно уметь после повторения}

\begin{itemize}
    \item определять ограничения, возникающие из-за знаменателей и квадратных корней;
    \item записывать ОДЗ как систему условий и находить их пересечение;
    \item решать квадратные неравенства методом интервалов;
    \item решать дробно-рациональные неравенства через нули числителя и знаменателя;
    \item правильно учитывать строгие и нестрогие знаки, а также кратность корней;
    \item отличать уравнения, которые используются для поиска критических точек, от исходного неравенства.
\end{itemize}

\topicsection{2. Базовые ограничения области определения}

\subsection*{2.1. Квадратный корень}

Если в выражении есть
\[
\sqrt{f(x)},
\]
то необходимо условие
\[
f(x)\ge 0.
\]

Корень может быть равен нулю, если сам корень не находится в знаменателе.

\subsection*{2.2. Знаменатель}

Если в выражении есть
\[
\frac{g(x)}{f(x)},
\]
то необходимо условие
\[
f(x)\ne 0.
\]

Числитель сам по себе ограничений не создаёт.

\subsection*{2.3. Корень в знаменателе}

Если в выражении есть
\[
\frac{g(x)}{\sqrt{f(x)}},
\]
то одновременно нужны условия
\[
f(x)\ge 0,
\qquad
\sqrt{f(x)}\ne 0.
\]

Их удобно объединить:
\[
\boxed{f(x)>0}.
\]

\subsection*{2.4. Несколько источников ограничений}

Все условия выполняются \key{одновременно}. Поэтому они записываются системой, а итоговая область определения является пересечением найденных множеств.

\[
\begin{cases}
f_1(x)\ge 0,\\
f_2(x)>0,\\
f_3(x)\ne 0.
\end{cases}
\]

\topicsection{3. Разобранный тип: корень в числителе и обычный знаменатель}

Рассмотрим функцию
\[
y=\frac{\sqrt{10+3x-x^2}}{x-3}.
\]

Ограничения:
\[
\begin{cases}
10+3x-x^2\ge 0,\\
x-3\ne 0.
\end{cases}
\]

Первое неравенство удобно привести к положительному старшему коэффициенту:
\[
10+3x-x^2\ge 0
\iff
x^2-3x-10\le 0.
\]

Корни соответствующего уравнения:
\[
x^2-3x-10=0,
\qquad
D=9+40=49,
\]
\[
x_1=-2,
\qquad
x_2=5.
\]

Следовательно,
\[
x\in[-2;5].
\]

Из-за знаменателя исключаем точку \(x=3\):
\[
\boxed{D(y)=[-2;3)\cup(3;5]}.
\]

\begin{center}
\begin{tikzpicture}[x=1.35cm,y=1cm]

\draw[-{Latex[length=2.5mm]},thick]
    (-3,0)--(6,0);

\foreach \x/\lab in {-2/-2,3/3,5/5}{
    \draw
        (\x,0.09)--(\x,-0.09)
        node[below=4pt] {$\lab$};
}

\draw[
    accent,
    line width=3pt
]
    (-2,0)--(3,0);

\draw[
    accent,
    line width=3pt
]
    (3,0)--(5,0);

\fill[accent]
    (-2,0)
    circle (2.2pt);

\draw[
    accent,
    fill=white,
    line width=1.2pt
]
    (3,0)
    circle (2.7pt);

\fill[accent]
    (5,0)
    circle (2.2pt);

\end{tikzpicture}
\end{center}

\subsection*{Ключевой вывод}

Ограничение на числитель появляется здесь только из-за квадратного корня. Сам факт нахождения выражения в числителе ограничения не создаёт.

% =========================================================
% Блок-схема 1
% =========================================================

\clearpage
\thispagestyle{empty}

\begin{center}
{\Large
\textcolor{darkaccent}{
\textbf{Алгоритм: как находить область определения}
}}
\end{center}

\vspace{0.9cm}

\begin{center}
\begin{tikzpicture}[node distance=1.25cm and 2.0cm]

\node[
    flowstart
]
(start)
{Начало};

\node[
    flowaction,
    below=of start
]
(scan)
{Просмотрите выражение и отметьте все знаменатели и квадратные корни};

\node[
    flowdecision,
    below=1.35cm of scan
]
(root)
{Есть выражение $\sqrt{f(x)}$?};

\node[
    flowbranch,
    below left=1.5cm and 0.2cm of root
]
(rootyes)
{Запишите $f(x)\ge 0$; если этот корень стоит в знаменателе, запишите $f(x)>0$};

\node[
    flowbranch,
    below right=1.5cm and 0.2cm of root
]
(rootno)
{Переходите к проверке знаменателей};

\node[
    flowdecision,
    below=2.1cm of root
]
(den)
{Есть знаменатель $g(x)$, ещё не учтённый условием на корень?};

\node[
    flowbranch,
    below left=1.55cm and 0.2cm of den
]
(denyes)
{Запишите $g(x)\ne 0$};

\node[
    flowbranch,
    below right=1.55cm and 0.2cm of den
]
(denno)
{Новых ограничений этого типа нет};

\node[
    flowaction,
    below=2.2cm of den
]
(system)
{Объедините все ограничения в систему и решите каждое из них};

\node[
    flowaction,
    below=of system
]
(intersect)
{Найдите пересечение полученных множеств};

\node[
    flowstart,
    below=of intersect
]
(finish)
{Запишите область определения};

\draw[flowarrow]
    (start)--(scan);

\draw[flowarrow]
    (scan)--(root);

\draw[flowarrow]
    (root)
    --node[flowlabel,left]{да}
    (rootyes);

\draw[flowarrow]
    (root)
    --node[flowlabel,right]{нет}
    (rootno);

\draw[flowarrow]
    (rootyes.south)
    --
    (den.north west);

\draw[flowarrow]
    (rootno.south)
    --
    (den.north east);

\draw[flowarrow]
    (den)
    --node[flowlabel,left]{да}
    (denyes);

\draw[flowarrow]
    (den)
    --node[flowlabel,right]{нет}
    (denno);

\draw[flowarrow]
    (denyes.south)
    --
    (system.north west);

\draw[flowarrow]
    (denno.south)
    --
    (system.north east);

\draw[flowarrow]
    (system)--(intersect);

\draw[flowarrow]
    (intersect)--(finish);

\end{tikzpicture}
\end{center}

\clearpage

% =========================================================
% Метод интервалов
% =========================================================

\topicsection{4. Квадратное неравенство и метод интервалов}

Для неравенства
\[
x^2-3x-10\le 0
\]
основная схема такова:

\begin{enumerate}
    \item заменить неравенство соответствующим уравнением
    \[
    x^2-3x-10=0;
    \]

    \item найти его корни:
    \[
    x=-2,
    \qquad
    x=5;
    \]

    \item отметить корни на числовой прямой;

    \item определить знак выражения хотя бы на одном интервале;

    \item при простых корнях чередовать знаки при переходе через каждый корень;

    \item выбрать интервалы, отвечающие знаку исходного неравенства.
\end{enumerate}

Для нестрогого неравенства корни включаются в ответ. Для строгого --- не включаются.

\subsection*{Удобная калибровка знаком}

Если \(0\) лежит на одном из интервалов, часто удобно подставить именно его. При этом достаточно определить \key{знак} значения, точное число обычно не требуется.

\topicsection{5. Дробно-рациональное неравенство}

Рассмотрим
\[
\frac{(x-3)(x-4)}
     {(x-1)(-x^2+7x-10)}
\ge 0.
\]

Разложим квадратный множитель знаменателя:
\[
-x^2+7x-10
=
-(x-2)(x-5).
\]

Критические точки получаем из \key{уравнений}:
\[
(x-3)(x-4)=0
\quad\Longrightarrow\quad
x=3,\ 4,
\]
\[
(x-1)(-x^2+7x-10)=0
\quad\Longrightarrow\quad
x=1,\ 2,\ 5.
\]

Точки \(1,2,5\) являются нулями знаменателя и обязательно выколоты. Точки \(3,4\) являются нулями числителя; при знаке \(\ge\) они включаются.

На крайнем левом интервале удобно проверить \(x=0\):
\[
\frac{(-3)(-4)}
     {(-1)(-10)}
=
\frac{12}{10}
>
0.
\]

Все критические точки здесь имеют нечётную кратность, поэтому знак меняется при переходе через каждую из них.

\begin{center}
\begin{tikzpicture}[x=1.45cm,y=1cm]

\draw[-{Latex[length=2.5mm]},thick]
    (0,0)--(6.2,0);

\foreach \x/\lab in {1/1,2/2,3/3,4/4,5/5}{
    \draw
        (\x,0.1)--(\x,-0.1)
        node[below=4pt] {$\lab$};
}

\node at (0.45,0.42) {$+$};
\node at (1.5,0.42) {$-$};
\node at (2.5,0.42) {$+$};
\node at (3.5,0.42) {$-$};
\node at (4.5,0.42) {$+$};
\node at (5.65,0.42) {$-$};

\draw[
    accent,
    fill=white,
    line width=1.2pt
]
    (1,0)
    circle (2.6pt);

\draw[
    accent,
    fill=white,
    line width=1.2pt
]
    (2,0)
    circle (2.6pt);

\fill[accent]
    (3,0)
    circle (2.2pt);

\fill[accent]
    (4,0)
    circle (2.2pt);

\draw[
    accent,
    fill=white,
    line width=1.2pt
]
    (5,0)
    circle (2.6pt);

\end{tikzpicture}
\end{center}

Отсюда
\[
\boxed{
x\in
(-\infty;1)
\cup
(2;3]
\cup
[4;5)
}.
\]

\topicsection{6. Кратность корня: когда знак не меняется}

Если множитель имеет вид
\[
(x-a)^2,
\]
то точка \(x=a\) имеет чётную кратность. При переходе через неё знак всего произведения или дроби \key{не меняется}.

Пример:
\[
\frac{(x-2)^2(x-1)}
     {x(x-3)(x+2)}
\ge 0.
\]

Критические точки:
\[
-2,\ 0,\ 1,\ 2,\ 3.
\]

Точка \(x=2\) --- корень чётной кратности, поэтому знак слева и справа от неё одинаков.

\subsection*{Типичная ошибка}

Если дискриминант квадратного уравнения равен нулю, корень существует:
\[
D=0
\quad\Longrightarrow\quad
x_1=x_2.
\]

Это один корень второй кратности, а не отсутствие корней.

% =========================================================
% Блок-схема 2
% =========================================================

\clearpage
\thispagestyle{empty}

\begin{center}
\begin{tikzpicture}[node distance=0.68cm]

\node[
    font=\Large\bfseries,
    text=darkaccent,
    align=center
]
(algtitle)
{Алгоритм: метод интервалов\\для дробного неравенства};

\node[
    flowstart,
    below=0.55cm of algtitle
]
(starttwo)
{Исходное неравенство};

\node[
    flowaction,
    below=of starttwo
]
(zero)
{Перенесите всё в одну сторону так, чтобы справа был $0$};

\node[
    flowaction,
    below=of zero
]
(factor)
{Разложите числитель и знаменатель на множители, если это возможно};

\node[
    flowaction,
    below=of factor
]
(equations)
{Решите уравнения: числитель $=0$ и знаменатель $=0$};

\node[
    flowaction,
    below=of equations
]
(points)
{Отметьте все критические точки на одной числовой прямой; нули знаменателя выколите};

\node[
    flowaction,
    below=of points
]
(sign)
{Определите знак на одном интервале удобной подстановкой};

\node[
    flowdecision,
    below=0.8cm of sign
]
(mult)
{Переходим через корень нечётной кратности?};

\node[
    flowbranch,
    below left=0.9cm and 0.2cm of mult
]
(change)
{Поменяйте знак};

\node[
    flowbranch,
    below right=0.9cm and 0.2cm of mult
]
(keep)
{Сохраните знак};

\node[
    flowaction,
    below=1.45cm of mult
]
(choose)
{Выберите интервалы нужного знака; включите допустимые нули числителя при $\le$ или $\ge$};

\node[
    flowstart,
    below=0.68cm of choose
]
(finishtwo)
{Запишите ответ};

\draw[flowarrow]
    (starttwo)--(zero);

\draw[flowarrow]
    (zero)--(factor);

\draw[flowarrow]
    (factor)--(equations);

\draw[flowarrow]
    (equations)--(points);

\draw[flowarrow]
    (points)--(sign);

\draw[flowarrow]
    (sign)--(mult);

\draw[flowarrow]
    (mult)
    --node[flowlabel,left]{да}
    (change);

\draw[flowarrow]
    (mult)
    --node[flowlabel,right]{нет}
    (keep);

\draw[flowarrow]
    (change.south)
    --
    (choose.north west);

\draw[flowarrow]
    (keep.south)
    --
    (choose.north east);

\draw[flowarrow]
    (choose)--(finishtwo);

\end{tikzpicture}
\end{center}

\clearpage

% =========================================================
% Итоговый чек-лист
% =========================================================

\topicsection{7. Что особенно важно не перепутать}

\begin{enumerate}
    \item
    \key{Для поиска критических точек решаются уравнения.}
    Исходное неравенство временно заменяется уравнениями только для нахождения нулей числителя и знаменателя.

    \item
    \key{Знаменатель не может быть равен нулю.}
    Его корни всегда исключаются из ответа.

    \item
    \key{Корень в числителе допускает ноль.}
    Условие имеет вид
    \[
    f(x)\ge 0.
    \]

    \item
    \key{Корень в знаменателе ноль не допускает.}
    Условие имеет вид
    \[
    f(x)>0.
    \]

    \item
    \key{При умножении или делении неравенства на отрицательное число знак меняется на противоположный.}

    \item
    \key{Чётная кратность критической точки не меняет знак.}
    Нечётная кратность меняет знак.

    \item
    \key{Система условий означает пересечение множеств.}
    После решения отдельных условий требуется собрать общий ответ.
\end{enumerate}

\topicsection{8. Мини-чек перед записью ответа}

Перед финальной строкой проверьте:

\begin{itemize}
    \item все ли знаменатели учтены;
    \item все ли подкоренные выражения учтены;
    \item правильно ли выбраны строгие или нестрогие условия;
    \item выколоты ли нули знаменателя;
    \item включены ли допустимые нули числителя при нестрогом знаке;
    \item учтена ли кратность каждого корня;
    \item выполнено ли пересечение всех условий ОДЗ.
\end{itemize}

% =========================================================
% Тренировочный блок
% =========================================================

\clearpage
\thispagestyle{plain}

\begin{center}
{\Large
\textcolor{darkaccent}{
\textbf{Тренировочный блок}
}
}\\[0.35em]

{\large
Путь А \textbullet{} по нарастающей сложности
}
\end{center}

\vspace{0.7cm}

\begin{enumerate}[label=\textbf{\arabic*)}]

    \item
    Найдите область определения функции
    \[
    y=
    \frac{\sqrt{x+4}}
         {x-2}.
    \]

    \item
    Найдите область определения функции
    \[
    y=
    \frac{\sqrt{x^2-x-6}}
         {x+1}.
    \]

    \item
    Решите неравенство
    \[
    x^2-5x+6\le 0.
    \]

    \item
    Решите неравенство
    \[
    \frac{(x-1)(x-4)}
         {(x+2)(x-3)}
    \ge 0.
    \]

    \item
    Решите неравенство
    \[
    \frac{(x-2)^2(x+1)}
         {x(x-4)(x+3)}
    \le 0.
    \]

\end{enumerate}

% =========================================================
% Ответы
% =========================================================

\clearpage
\thispagestyle{plain}

\begin{center}
{\Large
\textcolor{darkaccent}{
\textbf{Ответы к тренировочному блоку}
}}
\end{center}

\vspace{0.7cm}

\begin{table}[ht]
\centering
\caption{Краткие ответы}

\begin{tabularx}{\linewidth}{@{}cX@{}}
\toprule
\textbf{№} & \textbf{Ответ} \\
\midrule

1 &
$[-4;2)\cup(2;+\infty)$
\\

2 &
$(-\infty;-2]\cup[3;+\infty)$
\\

3 &
$[2;3]$
\\

4 &
$(-\infty;-2)\cup[1;3)\cup[4;+\infty)$
\\

5 &
$(-3;-1]\cup(0;4)$
\\

\bottomrule
\end{tabularx}
\end{table}

\vspace{0.9cm}

\textbf{Примечание к №5.}
Точка \(x=2\) имеет чётную кратность, поэтому знак при переходе через неё не меняется; она входит в интервал \((0;4)\).

\end{document}